\documentclass[10pt,a4paper]{amsart}
\usepackage[margin=19mm]{geometry}
\usepackage{amsmath,amssymb,mathtools}
\usepackage[scr=rsfs,cal=euler]{mathalfa}
\usepackage{xfrac}
\usepackage[unicode,hidelinks,bookmarks=false]{hyperref}
\newcommand{\ie}{i.e.\ }
\newcommand{\cf}{cf.\ }
\newcommand{\resp}{respectively\ }
\newcommand{\Subsection}{\subsection}
\providecommand{\nobreakdash}{\nobreak\hbox{-}}
\setlength{\emergencystretch}{2em}
\multlinegap0pt
\usepackage{xcolor}
\usepackage[shortlabels]{enumitem}
\usepackage{amssymb}
\usepackage{mathtools}
\newtheorem{definition}{Definition}
\newtheorem{theorem}{Theorem}
\newcommand{\CC}{\mathbb{C}}
\newcommand{\F}{\mathbb{F}}
\newcommand{\ZZ}{\mathbb{Z}}
\title{A Finiteness Theorem for Quartic K3-Fibred Calabi--Yau Threefolds in Scrolls}
\author{Geoffrey Mboya}
\date{Source: arXiv:2608.21934v1 (CC BY 4.0)}
\begin{document}
\maketitle

\noindent\emph{Source and licence.} A Finiteness Theorem for Quartic K3-Fibred Calabi--Yau Threefolds in Scrolls. Version arXiv:2608.21934v1 (CC BY 4.0), distributed by arXiv under the Creative Commons Attribution 4.0 International licence, http://creativecommons.org/licenses/by/4.0/, as recorded in the arXiv licence metadata for this exact version and shown on the abstract page https://arxiv.org/abs/2608.21934v1. Full licence notice: \texttt{licenses/arxiv-2608.21934-CC-BY-4.0.txt}.\par\smallskip

\noindent\textit{Editorial scope.} Counted unit: the article's theorem thm:system (source0.tex lines 197--202). The article's complete original proof of it is delivered with it (source0.tex lines 204--206, one \texttt{proof} environment of 3 lines). No mathematics is written, repaired or re-derived: the statement and the proof are the article's own lines, in the article's own notation, and no retained line is substituted or re-flowed.  Retained uncounted as the source-local context the proof needs: lines 166--172; lines 193--195.  Nothing was omitted from any retained span, so the package carries no omission record.  No retained line contains a citation, so the wrapper carries no bibliography and no bibliography entry.  No retained line contains a figure environment, an image-inclusion command or drawing code, so no image is delivered and the package declares no asset dependency.  Editorial formatting only, in addition: an AMS \texttt{amsart} preamble that restates the article's own declarations and macros used by the retained lines, namely the environments \texttt{definition}, \texttt{theorem} on separate counters, together with the standard packages those lines need.  The retained environments print in this document as Definition 1, Theorem 1 in order of appearance, whereas the article prints the same items with the numbering of its own full document.  The complete native source of the article as distributed in the arXiv e-print for this exact version is bundled unchanged as \texttt{sources/208287/source0.tex}.
	\begin{definition}\label{def:CY}
		An orbifold hypersurface $X\subset\F^m_n$ fibred by quartic K3 surfaces is a \emph{Calabi--Yau threefold} if $H^0(K_X)=\CC$, with $\beta\in\ZZ_n$ satisfying
		\[
		m_1q_1+\cdots+m_4q_4\equiv\beta\equiv m_1+\cdots+m_4\pmod n
		\]
		for $(q_j)\vdash4$.
	\end{definition}

	\begin{equation}\label{eq:canon}
		m_2+m_3+m_4\ge n+1.
	\end{equation}

	\begin{theorem}\label{thm:system}
		For $0=m_1\le m_2\le m_3\le m_4$ and $n\ge2$, $X\subset\mathbb{F}^m_n$ is Calabi--Yau with at worst canonical singularities along $C_1$ iff
		\begin{equation}\label{eq:system}
			m_2+m_3+m_4\equiv\beta\pmod n,\qquad m_2+m_3+m_4\ge n+1,\qquad 0\le m_2\le m_3\le m_4.
		\end{equation}
	\end{theorem}

	\begin{proof}
		By Definition \ref{def:CY}, $X$ is Calabi--Yau iff $\beta\equiv m_1+\cdots+m_4=m_2+m_3+m_4\pmod n$, the first condition of \eqref{eq:system}. By adjunction, canonicity of $X$ along $C_1$ is governed by the same age computation as the ambient orbifold, sharpened by one unit as in \eqref{eq:canon}, giving $m_2+m_3+m_4\ge n+1$; the ordering $0\le m_2\le m_3\le m_4$ is the normalisation of \S3.1. The same argument with indices permuted controls canonicity along $C_2,C_3,C_4$, and since $m_1=0$ makes $m_2+m_3+m_4$ the common value of the three remaining partial sums, condition \eqref{eq:canon} at $C_1$ already implies it at the other $C_j$.
	\end{proof}

\end{document}
